\documentclass[conference]{IEEEtran}
\IEEEoverridecommandlockouts

\usepackage{cite}
\usepackage{amsmath,amssymb,amsfonts}
\usepackage{algorithmic}
\usepackage{graphicx}
\usepackage{textcomp}
\usepackage{xcolor}

% Load dynamically generated numbers
% Auto-generated by build_numbers.py. Do not edit manually.
\newcommand{\CnnGruFone}{0.858}
\newcommand{\XgbFone}{0.868}
\newcommand{\CeilingDetect}{67.5\%}
\newcommand{\HiddenPct}{31.9\%}
\newcommand{\FerragRepAcc}{100.00\%}
\newcommand{\FerragCleanAcc}{97.66\%}
\newcommand{\FidTopFive}{0.605}


\begin{document}

\title{Auditing Shortcut Learning in IoT Intrusion Detection: Leakage, Detection Ceilings and Faithful Explanations}

\author{
\IEEEauthorblockN{1\textsuperscript{st} Given Name Surname}
\IEEEauthorblockA{\textit{dept. name of organization (of Aff.)} \\
\textit{name of organization (of Aff.)}\\
City, Country \\
email address or ORCID}
}

\maketitle

\begin{abstract}
Reported near-perfect accuracy on IoT/IIoT intrusion benchmarks is rarely accompanied by an audit of non-behavioural shortcuts. We present a leakage audit protocol and apply it to Edge-IIoTset and CICIoT2023, identifying representation and capture-level artifacts. We quantify the irreducible-error ceiling under a cleaned feature policy, showing that \HiddenPct\ of attack packets are indistinguishable from normal traffic. We demonstrate how reported accuracy changes once shortcuts are removed, observing a drop from \FerragRepAcc\ to \FerragCleanAcc\ in previous pipelines. Within this controlled setting we compare per-packet models against strong tabular baselines, evaluate SHAP fidelity via retraining-based removal, and report edge deployment trade-offs.
\end{abstract}

\begin{IEEEkeywords}
IoT, Intrusion Detection, Explainable AI, Shortcut Learning, Leakage
\end{IEEEkeywords}

\section{Introduction}
Intrusion detection systems (IDS) for IoT networks often report near-perfect accuracy (e.g., 99.99\% \cite{ferrag2022edge}), which is highly suspicious given the complexity of the data \cite{arp2022dos, kapoor2023leakage, geirhos2020shortcut}. This paper audits these scores to identify underlying artifacts and shortcuts rather than proposing yet another black-box model. Our contributions are:
1) A leakage audit protocol that found encoding shortcuts in Edge-IIoTset and CICIoT2023.
2) A measured empirical detection ceiling under a cleaned feature set.
3) A retraining-based test of explanation fidelity across explainers \cite{lundberg2017unified, ribeiro2016why}.
4) An honest benchmark comparing neural networks (CNN-GRU \cite{cho2014learning}) to tabular baselines (XGBoost \cite{chen2016xgboost}) with validation-only thresholds \cite{akiba2019optuna}.

\section{Related Work}
Table~\ref{tab:related} compares our work against recent literature using these datasets \cite{munilla2026xai, sharma2025self, wang2025shap, udurume2026shap}. We note that previous works often do not report a check for data leakage or evaluate explanation fidelity using retraining.

\begin{table*}[htbp]
\caption{Related Work Comparison}
\label{tab:related}
\centering
\begin{tabular}{|l|l|l|l|l|}
\hline
\textbf{Paper} & \textbf{Dataset(s)} & \textbf{Reported} & \textbf{Shortcut Check} & \textbf{Explanation Eval} \\
\hline
Ferrag 2022 \cite{ferrag2022edge} & Edge-IIoTset & 99.99\% binary & No & n/a \\
Munilla 2026 \cite{munilla2026xai} & Edge-IIoT, etc. & No accuracy reported & No (all 63 features) & Robustness under DeepFool \\
Sharma 2025 \cite{sharma2025self} & BoT-IoT, etc. & 97.9--99.6\% & No & SHAP--LIME alignment \\
Wang 2025 \cite{wang2025shap} & UNSW-NB15, etc. & up to 99.8\% & No & None \\
Udurume 2026 \cite{udurume2026shap} & UNSW-NB15, etc. & 97.5\% F1 & No & Zeroing at inference \\
\textbf{Ours} & Edge-IIoTset, CICIoT2023 \cite{neto2023ciciot} & \CnnGruFone\ (honest) & \textbf{Yes, two found} & \textbf{Retraining with random controls} \\
\hline
\end{tabular}
\end{table*}

\section{Datasets}
We utilize Edge-IIoTset and CICIoT2023 \cite{neto2023ciciot}. Edge-IIoTset suffers from file-level labels, while CICIoT2023 contains window-size artifacts which we successfully mitigated.

\section{Audit Protocol}
We implement a robust protocol including representation scans and identifier removals (e.g., \textit{tcp.seq}). We empirically demonstrated a detection ceiling, discovering that \HiddenPct\ of attack packets share identical feature patterns with normal packets, capping maximum detection to \CeilingDetect.

\section{Models and Training}
We evaluated CNN-GRU, MLP, 1D-CNN, GRU, and XGBoost. Hyperparameters were tuned via Optuna \cite{akiba2019optuna}, and thresholds were chosen exclusively on the validation set.

\section{Explanation and Fidelity}
We evaluated SHAP \cite{lundberg2017unified} using retraining fidelity. Removing the top-5 SHAP features dropped the Macro-F1 to \FidTopFive, confirming the model's reliance on these features, unlike zero-at-inference methods \cite{udurume2026shap}.

\section{Results}
\subsection{Shortcuts and the replication}
A standard pipeline on this dataset reaches near-perfect accuracy (\FerragRepAcc) through shortcuts. With every shortcut our audit detected removed, this drops to \FerragCleanAcc.
\subsection{Honest performance}
After cleaning, the CNN-GRU achieves an F1 of \CnnGruFone, and XGBoost slightly better at \XgbFone.
\subsection{Ceiling and label noise}
All networks reach the empirical ceiling under the cleaned feature set and this split.
\subsection{Fidelity}
Retraining without top features empirically demonstrated SHAP's faithfulness, dropping F1 to \FidTopFive.

\section{Discussion}
Near-perfect scores on these datasets often indicate artifact learning \cite{engelen2021trouble}. Dataset creators must label background traffic appropriately.

\section{Threats to Validity}
The evaluation relies on a random split where 99.5\% of test rows share a pattern with training. SHAP values are approximate. No physical hardware test has been conducted yet.

\section{Conclusion and Future Work}
Future work will explore sequences for the GRU and deployment on edge devices like Raspberry Pi.

\bibliographystyle{IEEEtran}
\bibliography{references}

\end{document}
